The broad scan evaluates primary reconstructions through asymmetric
distance computation: the query remains FP32 while a stored code selects
reconstructed coordinates \cite{pq,jhq}. For transformed query $q'$, the
ordinary table for primary subspace $m$ is
\begin{equation}
 T_m[z]=\|q'_m-C_m[z]\|_2^2,\qquad
 \primarydist(q,x)=\sum_{m=0}^{M-1}T_m[z_m(x)].
 \label{eq:5-primary-lut}
\end{equation}
With an eight-bit code, this requires 256 entries per subspace. The
opportunity here comes from the codebook's internal Cartesian structure,
not just the fact that the vector is partitioned into PQ subspaces.

\subsection{Factorisation without changing the quantiser}

For the tested $D_s=8$ Cartesian codebook, the high nibble $h$ selects
the first four reconstructed coordinates and the low nibble $l$ selects
the remaining four. Thus $C_m[16h+l]=(A_m[h],B_m[l])$, with no
cross-dependence between the two coordinate groups. Define
\begin{equation}
 T_m^{\rm hi}[h]=\|q'_{m,\rm hi}-A_m[h]\|_2^2,\qquad
 T_m^{\rm lo}[l]=\|q'_{m,\rm lo}-B_m[l]\|_2^2.
 \label{eq:5-lut-halves}
\end{equation}
Squared-L2 additivity over disjoint coordinates gives
\begin{equation}
 T_m[z]=T_m^{\rm hi}[z\!\gg\!4]
        +T_m^{\rm lo}[z\mathbin{\&}15].
 \label{eq:5-lut-factorisation}
\end{equation}
Only 32 entries are stored, rather than 256. This preserves the
reconstruction and the real-arithmetic score; no code retraining,
extra quantisation or change in the candidate universe is involved.

The builder reads representative codewords $16h$ and $l$ from the
existing full codebook and evaluates only their corresponding coordinate
slices. The zero digits in the other slice are irrelevant. This is valid
only when all codewords sharing a nibble share that slice. The
\texttt{verify\_split\_lut} check compares the stored codeword slices
against their representatives and rejects incompatible codebooks.
An arbitrary learned 256-centroid PQ subcodebook need not satisfy this
condition: dividing its code labels into nibbles does not factor its
centroid coordinates.

The mathematical identity does not promise bitwise FP32 equality. A
full eight-coordinate accumulation and two four-coordinate sums followed
by an addition associate the operations differently. The implementations
therefore retain FP32 tables and separately check numerical and ranking
effects. Replacing table entries with FP16 would introduce another
rounding operation and is a different design choice.

\subsection{Grouping, construction work and resource cost}

The v71 ablation generalises the split to $g=\glut\in\{1,2,4,8\}$
contiguous coordinate groups whose bit fields match the Cartesian
digits. For eight bits and $D_s=8$, each group has $2^{8/g}$ entries;
the total entries and FP32 storage per query are
\begin{equation}
 E(g)=g2^{8/g},\qquad W_{\rm LUT}(g)=4ME(g)\ \text{bytes}.
 \label{eq:5-lut-cost}
\end{equation}
Independent entry construction evaluates $D_s2^{8/g}$ coordinate
contributions per subspace. Two groups therefore reduce table storage
eightfold and coordinate work sixteenfold relative to independently
constructing every full-table entry. These are table-level cost ratios,
not predicted application speedups.

\begin{table}[tbp]
\centering
\caption{Cartesian table trade-offs for an eight-bit, eight-coordinate
primary subspace. Combining additions are the algebraic minimum for
joining the looked-up values, excluding accumulation across subspaces;
they are not measured instruction counts.}
\label{tab:5-lut-resources}
\small
\begin{tabular}{@{}rrrrrr@{}}
\toprule
$\glut$ & Entries & FP32 bytes & Lookups & Combining adds & Build terms \\
\midrule
1 & 256 & 1,024 & 1 & 0 & 2,048 \\
2 & 32 & 128 & 2 & 1 & 128 \\
4 & 16 & 64 & 4 & 3 & 32 \\
8 & 16 & 64 & 8 & 7 & 16 \\
\bottomrule
\end{tabular}
\end{table}

The table is first built in global workspace. A scan block stages it
cooperatively in shared memory when the table and that kernel's candidate
buffer fit the device's configured limit; a block barrier completes
staging before lookup. Otherwise it reads the global table. The shared
footprint must include selection scratch: $W_{\rm LUT}$ alone does not
determine residency. Grouped scan blocks may each stage the same query's
table, coupling this design to \Cref{sec:5-5}.

At $M=96$ and 384, full FP32 tables require 96 and 384\,KiB per query;
two-group tables require 12 and 48\,KiB. The smaller footprint creates
placement opportunities, but four and eight groups make more lookups and
combine more values even though both store only 16 entries per subspace.
Index extraction, dependent lookups, synchronisation and residency all
contribute to the resulting trade-off. Table size alone cannot identify
the fastest design or establish a unique cache/bank-conflict explanation.

Small-table execution has substantial precedent in PQ Fast Scan and
Quicker ADC \cite{fastscan,quickeradc}. This implementation applies exact
coordinate factorisation to JHQ's existing Cartesian codebook and measures
the GPU grouping trade-off; it does not claim to invent small LUTs.
\Cref{sec:7-3} reports the controlled v71 sweep: two groups lead in
22/24 cells, with the other two differences below 1\%. That result supports
the selected operating point rather than a universal optimum.
